\documentclass[10pt]{article}
\usepackage[margin=0.75in]{geometry}
\usepackage[T1]{fontenc}
\usepackage{lmodern}
\usepackage{microtype}
\usepackage{amsmath}
\usepackage{graphicx}
\usepackage{hyperref}
\hypersetup{colorlinks=true,urlcolor=blue,linkcolor=blue}
\setlength{\parindent}{0pt}
\setlength{\parskip}{0.35em}
\sloppy

\title{\texorpdfstring{\(D^+ \to K\pi\pi\) Mass Peak Search}{D+ to Kpipi Mass Peak Search}}
\author{Evan Zhang}
\date{Codex export updated 2026-06-23}

\begin{document}
\maketitle
\vspace{-1.2em}

\section*{Goal}

This is a separate task from the \(D^0\) analysis-note reproduction. The
immediate goal is only to obtain a credible \(D^+ \to K\pi\pi\) mass peak from
2023 PbPb UPC data. Final \(D^+\) cut derivation, efficiencies, normalization,
and cross-section extraction are not attempted here yet.

\section*{Current Workflow}

\begin{enumerate}
\item Start from the same 2023 HIForward0 MINIAOD input family used for the
\(D^0\) reproduction.
\item Run the 2023 forest config with Dfinder configured for
\(D^+/D^- \to K\pi\pi\):
\[
\texttt{Dchannel} =
(0,0,1,1,0,0,0,0,0,0,0,0,0,0,0,0,0,0).
\]
The full 18-entry vector is required because Dfinder checks that every cut
vector has the same channel length.
\item Keep a minimal D+ forest containing event bookkeeping, track/PF/ZDC
diagnostics, and \texttt{Dfinder/ntDkpipi}.
\item Plot \(\texttt{Dmass}\) with transparent diagnostic selections:
raw Dfinder \(K\pi\pi\), a basic kinematic window, and a loose topology
window.
\item Scale with CRAB only after bounded local and tiny-CRAB smokes validate
the config.
\end{enumerate}

\section*{Differences from the \(D^0\) AN Workflow}

This document is not a replacement for the \(D^0\) analysis-note reproduction.
It is a first \(D^+\) peak-finding workflow built on the same 2023 UPC input
data family. The main differences are:

\begin{itemize}
\item The reconstructed channel changes from
\[
D^0/\overline{D^0}\to K\pi
\]
to
\[
D^+/D^-\to K\pi\pi .
\]
The plotted tree is therefore \texttt{Dfinder/ntDkpipi}, not
\texttt{Dfinder/ntDkpi}.
\item The Dfinder channel mask changes from the \(D^0\)-enabled channels to the
18-entry vector with only the \(D^+/D^-\) \(K\pi\pi\) channels enabled:
\[
\texttt{Dchannel} =
(0,0,1,1,0,0,0,0,0,0,0,0,0,0,0,0,0,0).
\]
\item The candidate mass is the \(K\pi\pi\) invariant mass under Dfinder's
kaon/pion mass hypotheses. There is no explicit kaon/pion PID in this
diagnostic workflow.
\item The \(D^0\) AN topological cuts are not reused as final \(D^+\) cuts.
The current loose/medium/tight topology selections are diagnostic peak-finding
cuts only. A real \(D^+\) analysis would need its own optimization, sideband
checks, signal efficiency, and mass-shape validation.
\item The \(D^0\) AN corrections, yields, efficiencies, systematic
uncertainties, and cross-section extraction are not applied here. This note
only demonstrates that the recreated foresting path can produce a visible
\(D^+ \to K\pi\pi\) mass peak.
\item The event-level UPC data source and foresting machinery are inherited
from the \(D^0\) reproduction infrastructure, but the output forests are newly
produced D+ forests rather than the \(D^0\) selected skim used in the AN.
\end{itemize}

\section*{Smoke Validation}

The first fixed local smoke used 100 MINIAOD events. The earlier zero-candidate
attempt was traced to a malformed 14-entry \texttt{Dchannel} override; after
restoring the required 18 entries, Dfinder stopped printing the size-mismatch
warning and populated \texttt{ntDkpipi}.

\begin{itemize}
\item No-superfilter smoke: 100 \texttt{ntDkpipi} event entries.
\item Normal UPC-filtered smoke: 11 filtered \texttt{ntDkpipi} event entries.
\item Candidate count before diagnostic mass/kinematic cuts: 8,130.
\item Diagnostic mass plot: 8,111 raw Kpipi candidates, 1,284 basic kinematic
candidates, and 34 loose-topology candidates.
\end{itemize}

\begin{figure}[h]
\centering
\includegraphics[width=0.82\linewidth]{support/figures/dplus_peak/dplus_mass_smoke100.png}
\caption{Smoke diagnostic from 100 input events. This validates D+ branch
production and plotting only; it is not enough statistics to claim a
\(D^+\) peak.}
\end{figure}

\section*{CRAB Status and Medium Run}

A one-job CRAB smoke was first attempted with the physically bounded
100-event MINIAOD file, rather than relying on CRAB-side \texttt{maxEvents}.
That task failed before physics processing because the tiny file is on
CERNBOX under \texttt{/eos/user/e/evzhang/...}, while CRAB resolved the
submitted \texttt{/store/user/...} LFN through the CMS \texttt{/eos/cms}
namespace. The failed tiny task was canceled to prevent retries.

A replacement one-file CRAB smoke was submitted with a real HIForward0 MINIAOD
LFN. It produced output successfully. The automated chain then submitted the
100-file medium run.

\begin{itemize}
\item Canceled tiny task:
\texttt{260622\_232410:evzhang\_crab\_dplus\_minimal\_tiny\_crab\_smoke\_20260622\_2323}.
\item Completed one-file smoke task:
\texttt{260622\_233514:evzhang\_crab\_dplus\_onefile\_crab\_smoke\_20260622\_2338}.
\item One-file smoke project directory:
\begin{quote}\footnotesize
\path{/afs/cern.ch/work/e/evzhang/codex-work/analysis/dplus-peak/crab-smoke/crab-projects/crab_dplus_onefile_crab_smoke_20260622_2338}
\end{quote}
\item Medium task:
\texttt{260623\_062056:evzhang\_crab\_dplus\_medium\_100files\_20260622\_overnight}.
\item Medium project directory:
\begin{quote}\footnotesize
\path{/afs/cern.ch/work/e/evzhang/codex-work/analysis/dplus-peak/crab-smoke/crab-projects/crab_dplus_medium_100files_20260622_overnight}
\end{quote}
\item Active chain monitor:
\texttt{research/cms-dplus-analysis/scripts/launch\_dplus\_medium\_after\_smoke.sh}.
\item Chain behavior: monitor for 100 output ROOT files and run the D+ mass
plotter when they are all visible.
\end{itemize}

\section*{Preliminary Medium-Run Checkpoint}

A checkpoint plot was made while the medium CRAB task was still running. To
avoid scanning all visible multi-GB forest outputs interactively, the checkpoint
used the first four visible medium-run ROOT files.

\begin{itemize}
\item Checkpoint label:
\texttt{dplus\_medium\_prelim4\_readable\_20260623T1100Z}.
\item Visible medium outputs at checkpoint time: 34.
\item Used outputs: 4.
\item Candidates in the four-file checkpoint:
615,362 loose topology, 9,072 medium topology, and 202 tight topology.
\item Medium-topology diagnostic fit:
\[
m = 1.87381~\mathrm{GeV},\qquad
\sigma = 0.00649~\mathrm{GeV}.
\]
\end{itemize}

This was close to the nominal \(D^+\) mass, but the four-file plot was still
dominated by smooth combinatorial background and statistical fluctuations.

A stronger checkpoint was then run over 20 visible medium-run ROOT files. The
medium topology view was still background dominated, but the tight topology
view produced a visible peak near the nominal \(D^+\) mass.

\begin{itemize}
\item Checkpoint label:
\texttt{dplus\_medium\_prelim20\_toposcan\_20260623T1105Z}.
\item Visible medium outputs at checkpoint time: 39.
\item Used outputs: 20.
\item Candidates in the 20-file checkpoint:
2,680,780 loose topology, 38,122 medium topology, and 932 tight topology.
\item Tight-topology diagnostic fit:
\[
m = 1.87190~\mathrm{GeV},\qquad
\sigma = 0.01486~\mathrm{GeV}.
\]
\end{itemize}

The honest current classification is therefore \emph{credible preliminary
peak}. It is not a final analysis result: these are diagnostic topology cuts,
not yet final optimized or efficiency-validated \(D^+\) selections.

At the user's request, the same plotter was rerun over all 59 medium-run ROOT
outputs visible at that time. This is the current best-statistics checkpoint.

\begin{itemize}
\item Checkpoint label:
\texttt{dplus\_medium\_prelim59\_toposcan\_20260623T1230Z}.
\item Visible and used medium outputs: 59.
\item Candidates in the 59-file checkpoint:
8,497,790 loose topology, 118,117 medium topology, and 2,570 tight topology.
\item Medium-topology diagnostic fit:
\[
m = 1.87164~\mathrm{GeV},\qquad
\sigma = 0.00803~\mathrm{GeV}.
\]
\item Tight-topology diagnostic fit:
\[
m = 1.87153~\mathrm{GeV},\qquad
\sigma = 0.01142~\mathrm{GeV}.
\]
\end{itemize}

The 59-file tight-topology plot is the clearest current D+ peak result.

\begin{figure}[h]
\centering
\includegraphics[width=0.82\linewidth]{support/figures/dplus_peak/dplus_mass_dplus_medium_prelim4_readable_20260623T1100Z_medium_topology.png}
\caption{Four-file medium-run checkpoint under the medium topology diagnostic.
The fitted mean is near the expected \(D^+\) mass, but the shape is not yet a
robust visible peak.}
\end{figure}

\begin{figure}[h]
\centering
\includegraphics[width=0.82\linewidth]{support/figures/dplus_peak/dplus_mass_dplus_medium_prelim20_toposcan_20260623T1105Z_tight_topology.png}
\caption{Twenty-file medium-run checkpoint under the tight topology diagnostic.
This is the first credible preliminary \(D^+ \to K\pi\pi\) peak from the
recreated D+ medium forests.}
\end{figure}

\begin{figure}[h]
\centering
\includegraphics[width=0.82\linewidth]{support/figures/dplus_peak/dplus_mass_dplus_medium_prelim59_toposcan_20260623T1230Z_tight_topology.png}
\caption{Fifty-nine-file medium-run checkpoint under the tight topology
diagnostic. This is the current best-statistics preliminary \(D^+ \to K\pi\pi\)
peak from the recreated D+ forests.}
\end{figure}

\section*{Next Scale-Up}

The medium CRAB monitor now ignores the one failed job and waits for 99 visible
outputs. Once those outputs are present, it should automatically rerun the same
D+ mass plotter on the larger sample.

\section*{Durable Scripts}

\begin{itemize}
\item Local forest smoke:
\texttt{research/cms-dplus-analysis/scripts/run\_2023\_dplus\_forest\_smoke.sh}.
\item CRAB forest runner:
\texttt{research/cms-dplus-analysis/scripts/run\_2023\_dplus\_forest\_crab\_smoke.sh}.
\item Mass diagnostic macro:
\texttt{research/cms-dplus-analysis/scripts/plot\_dplus\_mass.C}.
\item Visible-output checkpoint helper:
\texttt{research/cms-dplus-analysis/scripts/plot\_visible\_dplus\_crab\_outputs.sh}.
\end{itemize}

\end{document}
